\section{Билет 62. Тепловая машина. КПД тепловой машины. Вычисление КПД цикла Отто (двигатель внутреннего сгорания) и сравнение с КПД цикла Карно.}

\begin{definition}
\textbf{Тепловой машиной} называется устройство, преобразующее внутреннюю энергию топлива в механическую работу с помощью кругового процесса. \textbf{КПД} машины определяется как $\eta = A/Q_1 = 1 - |Q_2|/Q_1$.
\end{definition}

\begin{theorem}[Цикл Отто и его КПД]
\textbf{Цикл Отто} --- идеализированный термодинамический цикл бензинового двигателя внутреннего сгорания. Состоит из четырёх процессов:
\begin{enumerate}
    \item 1--2: Адиабатическое сжатие рабочей смеси;
    \item 2--3: Изохорное подведение теплоты $Q_1$ (сгорание топлива);
    \item 3--4: Адиабатическое расширение (рабочий ход);
    \item 4--1: Изохорное отведение теплоты $|Q_2|$ (выпуск отработавших газов).
\end{enumerate}
Для изохорных процессов количество теплоты:
\begin{equation}
    Q_1 = \nu C_V (T_3 - T_2), \quad |Q_2| = \nu C_V (T_4 - T_1).
\end{equation}
КПД цикла выражается через температуры:
\begin{equation}
    \eta_{\text{Отто}} = 1 - \frac{T_4 - T_1}{T_3 - T_2}. \label{eq:otto_eff_temp}
\end{equation}
Для адиабат 1--2 и 3--4 справедливы соотношения $T V^{\gamma-1} = \text{const}$:
\begin{equation}
    T_1 V_1^{\gamma-1} = T_2 V_2^{\gamma-1}, \quad T_4 V_1^{\gamma-1} = T_3 V_2^{\gamma-1}.
\end{equation}
Отсюда следует:
\begin{equation}
    \frac{T_2}{T_1} = \frac{T_3}{T_4} = \left(\frac{V_1}{V_2}\right)^{\gamma-1} = n^{\gamma-1},
\end{equation}
где $n = V_1/V_2$ --- \textbf{степень сжатия}. Тогда $\frac{T_4}{T_3} = \frac{T_1}{T_2} = \frac{1}{n^{\gamma-1}}$, и выражение \eqref{eq:otto_eff_temp} преобразуется к виду:
\begin{equation}
    \eta_{\text{Отто}} = 1 - \frac{1}{n^{\gamma-1}}. \label{eq:otto_eff_final}
\end{equation}
\end{theorem}

\begin{property}[Сравнение с циклом Карно]
Цикл Отто работает между максимальной температурой $T_3$ и минимальной $T_1$. КПД цикла Карно для этих же предельных температур равен:
\begin{equation}
    \eta_{\text{К}} = 1 - \frac{T_1}{T_3}.
\end{equation}
Сравнение показывает, что $\eta_{\text{Отто}} < \eta_{\text{К}}$. Это обусловлено тем, что в цикле Отто тепло подводится и отводится не изотермически, а изохорно. Средняя температура подвода тепла оказывается ниже $T_3$, а средняя температура отвода выше $T_1$. Согласно теореме Карно, любое отклонение от идеального изотермического теплообмена снижает КПД.
\end{property}

\begin{remark}
\textbf{Практические аспекты:}
\begin{itemize}
    \item КПД цикла Отто зависит только от степени сжатия $n$ и показателя адиабаты $\gamma$. Для бензиновых двигателей $n \approx 8 \text{--} 12$, $\gamma \approx 1.4$, что даёт теоретический $\eta_{\text{Отто}} \approx 0.56 \text{--} 0.63$. Реальный КПД составляет $25 \text{--} 35\%$ из-за трения, потерь тепла в стенки цилиндра, неполного сгорания и насосных потерь.
    \item Увеличение $n$ повышает КПД, но ограничено детонацией (самовоспламенением топливной смеси). Использование высокооктанового топлива и систем непосредственного впрыска позволяет безопасно увеличивать степень сжатия.
    \item Цикл Отто является частным случаем более общего термодинамического цикла, но его аналитическая простота и наглядность делают его стандартом для теоретического анализа поршневых ДВС.
\end{itemize}
\end{remark}
